% ============================================================
% base/templates/components.tex   内置组件（基线 · 仓库里的源）
%
% 页面文件只允许使用本文件里的命令 + 纯文本。
% 出现 \fontsize / \vspace / \setbeamercolor 就是设计债，
% 应该加组件，而不是在页面里写样式。
%
% 【这份文件不参与 deck 编译】
%   新建 deck 时只复制 `style.tex`、`main.tex`、`pages/`。
%   本文件与 format.tex 留在 skill 里，是给 deck**抄**的源：
%   deck 的 style.tex 自带等价内容，并另加载本 deck 的 .slide-forge/。
%
% 【要新增组件写到哪】
%   <deck>/slide/.slide-forge/templates/components.tex   ← 自动生成，不用问用户
%   evolved/templates/components.tex                     ← 吸收，用户点头后才做
%   本文件：只有用户明确同意"升入基线"时才动。
% ============================================================

% ============================================================
% 页面骨架
% ============================================================

% --- 有标题页：标题由 main.tex 给出，正文来自页面文件 ---
\newcommand{\DeckSlide}[2]{%
  \begin{frame}[t]{#1}%
    \vspace*{\SlideBodyTopSkip}%
    \input{#2}%
  \end{frame}%
}

% --- 无标题页：封面与致谢页用 ---
% #1 内容文件，#2 垂直对齐（c 居中 / t 顶对齐）
%
% 这两页不要页标题：封面本来就该是整页一块，致谢页只有两个字，
% 再顶一个页标题会显得多余、也会和正文页混淆。
% 做法是在这一页把 frametitle 模板清空（模板是局部的，不影响其他页），
% 内容用 \vfill 上下撑开，从而垂直居中。
\newcommand{\DeckBare}[2]{%
  \begin{frame}[#2]%
    \setbeamertemplate{frametitle}{}%
    \input{#1}%
  \end{frame}%
}

% ============================================================
% 正文行与分组
% ============================================================

% --- 正文行：全页只有这一种行 ---
% 姓名、单位、每一条要点，全部用同一个命令，因此字号、行距、
% 字重、颜色必然完全一致 —— 从代码层面就不可能出现不一致。
% 不要为"姓名"或"要点"另建命令：两个命令意味着两套样式，
% 迟早会写出不一样的效果。一个命令，一个样式。
\newcommand{\DeckLine}[1]{%
  \par
  \hangindent=1.1em\hangafter=1
  \noindent #1\par
  \vspace{0.30cm}%
}

% --- 分组留白 ---
% 唯一用来分组的手段：在两组信息之间插入一次更大的空白。
% 不用分隔线、不用小标题、不用字号差异。
\newcommand{\DeckGap}{%
  \par\vspace*{0.40cm}%
}

% --- 带时间标签的行 ---
% #1 时间标签，#2 内容
%
% 不用空格去凑对齐：标签字数不等（"大二寒假" 4 字、"本学期" 3 字），
% 手打空格永远对不齐，而且全角/半角混用会更乱。
% 这里用一个固定宽度的盒子把标签撑到统一宽度，内容自然左对齐。
\newlength{\DeckTagWidth}
\setlength{\DeckTagWidth}{5.6em}   % 约容纳 5 个全角字符
\newcommand{\DeckTagged}[2]{%
  \DeckLine{\makebox[\DeckTagWidth][l]{#1}#2}%
}

% ============================================================
% 条目（带链接）
% ============================================================

% --- 条目：一条两行（标题行 + 说明行）---
% #1 标题行（可内嵌 \DeckLink 让标题本身可点），#2 说明行（角色 / 出处）
% 两行之间比 \DeckLine 更紧（0.12cm），让它们读作一条完整记录；
% 条目之间用 0.34cm 拉开。标题折行时第二行与首行文字左对齐。
\newcommand{\DeckCite}[2]{%
  \par
  \hangindent=1.1em\hangafter=1
  \noindent #1\par
  \vspace{0.12cm}%
  \noindent #2\par
  \vspace{0.34cm}%
}

% --- 超链接 ---
% beamer 已加载 hyperref，直接用 \href。
% 颜色用强调色，不加下划线边框（hyperref 的默认方框很难看）。
\newcommand{\DeckLink}[2]{%
  \href{#1}{\color{SlideAccent}#2}%
}

% ============================================================
% 表格
% ============================================================

% --- 表格行距（可调旋钮）---
% 默认值 = 正文行距，因此表格与正文看起来是一套节奏。
% 某页需要更紧或更松时，在该页文件开头改这一个长度即可：
%     \setlength{\DeckTableLeading}{14pt}
% 注意：字号不在这里调。字号只有两档，全篇一致，不允许逐页改动。
\newlength{\DeckTableLeading}
\setlength{\DeckTableLeading}{\baselineskip}

% --- 数值表：「名称占满 + 数值右对齐」---
% #1 表体：名称 & 数值 \\
% 表宽自动等于 \linewidth。表内文字用 \SlideTextFont，字号与正文一致；
% 层次只靠加粗（列头）与列间距。不画竖线、不画外框。
\newcommand{\DeckCourseTable}[1]{%
  \par\noindent
  {\SlideTextFont
   \setlength{\baselineskip}{\DeckTableLeading}%
   \setlength{\tabcolsep}{0pt}%
   \renewcommand{\arraystretch}{1.0}%
   \noindent
   \begin{tabularx}{\linewidth}{@{}>{\raggedright\arraybackslash}X@{\hspace{0.16cm}}r@{}}
     #1
   \end{tabularx}}%
  \par
}

% --- 列表式两列表：「窄标签左 + 长文字占满左对齐」---
% 与 \DeckCourseTable 的区别只在列宽分配方向：
%   数值表是「名称占满、数值右对齐」；
%   本表是「短标签在左（窄）、长文字在右（占满、左对齐）」。
% 年份 + 描述、时间线这类内容用这个，
% 否则描述会被右对齐得零碎，且长条目会溢出列宽。
\newcommand{\DeckListTable}[1]{%
  \par\noindent
  {\SlideTextFont
   \setlength{\baselineskip}{\DeckTableLeading}%
   \setlength{\tabcolsep}{0pt}%
   \renewcommand{\arraystretch}{1.0}%
   \noindent
   \begin{tabularx}{\linewidth}{@{}l@{\hspace{0.55cm}}>{\raggedright\arraybackslash}X@{}}
     #1
   \end{tabularx}}%
  \par
}

% --- 表头与细线 ---
% 表头行：与正文同字号，只靠加粗区分
\newcommand{\DeckCourseHead}[2]{%
  \textbf{#1} & \textbf{#2} \\[0.04cm]
}
% 列表式表头
\newcommand{\DeckListHead}[2]{%
  \textbf{#1} & \textbf{#2} \\[0.04cm]
}
% 列头与第一行数据之间的细线。
% 用 \noalign 单独发 \hrule，避免 \hline\noalign 给表格带来的异常深度
% （实测那一手会让 5 行表的深度变成 77pt）。
\newcommand{\DeckCourseRule}{%
  \noalign{\vspace{0.09cm}\hrule height 0.4pt\vspace{0.13cm}}%
}

% ============================================================
% 图片
% ============================================================

% --- 两图并排 ---
% #1 左图路径，#2 右图路径，#3 左图高度，#4 右图宽度
%
% 不用 beamer 的 columns：它的栏间距会与给定宽度叠加，很难算准，
% 几次都报 Overfull。改成两图放在同一行、用 \hfill 撑开，
% 宽度完全由 \includegraphics 的 height/width 决定，可控。
%
% \includegraphics 的 width/height 是"缩放后"尺寸：只给 height 时，
% 宽度按原图比例推算，所以两张图若比例不同、都想等高，
% 就必须给其中一张显式 width（调用方算好传入）。
%
% 关键约束：两图宽度之和必须明显小于正文宽（14.25cm），留 0.3cm 以上。
% 只差 0.03cm 都会让第二张图被挤到下一行，并报 overfull vbox。
\newcommand{\DeckFigures}[4]{%
  \par\noindent
  \includegraphics[height=#3]{#1}%
  \hfill
  \includegraphics[width=#4]{#2}%
  \par
}
